\begin{python0}
from solutions import *; clear()
\end{python0}
\sectionthree{reinterpret\_cast and void *}

Recall that while you can do C-style casting for pointer types:

\begin{console}
int i = 42;
float * p = (float *)(&i);
\end{console}

Using C++ typecasting, the above is done using the \texttt{reinterpret\_cast}. Try this:

\begin{console}
int i = 42;
float * p = reinterpret_cast< float * >(&i);
\end{console}

You should try to do the above with \texttt{\EMPHASIZE{static\_cast}} and \texttt{\EMPHASIZE{const\_cast}} to check that neither of them works.

You should try this:

\begin{console}
#include <iostream>

int main()
{   
    int i = 42;
    float * p = reinterpret_cast< float * >(&i);
    *p = 0.0f;
    std::cout << i << '\n'; // you get 0 because the
                            // bit pattern of 0.0f
                            // happens to be the same
                            // bit pattern of 0.
    *p = 3.14f;
    std::cout << i << '\n'; // you do NOT get 3
                            // because the bit
                            // pattern of 3.14f is
                            // NOT 3.
    return 0;
}
\end{console}

The second output is \texttt{1078523331}. Why? Because the bit pattern for \texttt{3.14f} happens to be the bit pattern of the integer \texttt{1078523331}.

In general you should probably not use \texttt{reinterpret\_cast} of pointers until you have taken CISS360 where you will learn more about bit representations and C/C++ bit programming.

Note that in the above \texttt{p} points to a \texttt{float} (the smaller version of \texttt{double}). So C++ will treat \texttt{*p} is a \texttt{float}. Sometimes (in the future), you might simply \EMPHASIZE{want an address} and you \EMPHASIZE{don't really want to interpret the value at that address} as an \texttt{int} or \texttt{double} or \texttt{float} or \texttt{char} or \texttt{bool} or anything at all. In that case, the pointer type you want is called

\begin{center}
\EMPHASIZE{void *}
\end{center}

Of course if \texttt{q} is a \texttt{void *}, you can print the address
stored in \texttt{q}, you can give that address value to another pointer
variable. But you should not work with \texttt{*q} since C++ does not know
how to interpret the bits at that address. Run this:

\begin{console}
int i = 42;
float * p = reinterpret_cast< float * >(&i);
*p = 0.0f;
std::cout << i << '\n';
*p = 3.14f;
std::cout << i << '\n';

void * q = reinterpret_cast< float * >(&i);
std::cout << q << '\n';    // OK
std::cout << (*q) << '\n'; // ARGH!!!! WRONG!!!
\end{console}

Run this too where you store an \texttt{int*} value in a \texttt{void*} pointer, pass the value of the pointer to another \texttt{void*} pointer, cast the \texttt{void*} value to a \texttt{double*}:

\begin{console}
int * p0 = new int;
void * p1 = p0;
void * p2 = reinterpret_cast< void * >(p0);
double * p3 = reinterpret_cast< double * >(p2);
std::cout << p0 << ' ' << p1 << ' ' << p2 << ' '
          << p3 << '\n';
\end{console}

As I mentioned earlier, in the area of OpenCV and OpenGL, you'll see
that many software engineers do not use the C++ style casting. For
instance one function in OpenGL is
\begin{center}
\verb!void glVertexAttribPointer(GLuint index,                                                                   GLint size,                                                                     GLenum type,                                                                    GLboolean normalized,                                                           GLsizei stride,                                                                 const void * pointer);!
\end{center}
and you'll see that most OpenGL programmers will simply use the simple
C/C++ cast like this to call the above function:
\begin{center}
\verb!glVertexAttribPointer(0,                                                                              3,                                                                       GL_FLOAT,                                                                       GL_FALSE,                                                           5 * sizeof (GLfloat),                                                                   (GLvoid*) 0);!
\end{center}

(note: GLvoid is just a typedef of void.)

